\section{Limitations and Threats to Validity}
\label{sec:limitations}

We state the boundaries of the evidence plainly; several are active work items
rather than permanent limits.

\paragraph{Discovery vs.\ confirmation.} The precision result
(\S\ref{sec:olmo}(c)) is at present a configuration \emph{discovery} on a small
oracle-scored sample plus a deductive source-level argument, partially corroborated
by two ordinary-path runs (below); it has not been confirmed on \emph{held-out}
instances. Matching a small
sample does not uniquely identify a producing deployment, and the deductive leg is
conditional on the unrecorded \texttt{transformers} version.
{\color{orange}\textsf{\small[\textsc{pending}: held-out OLMo confirmation]}}

\paragraph{The confirmation experiment: begun, not completed to protocol.}
\label{sec:olmo-confirm}
Two ordinary-path \texttt{ifeval} runs at float32 exist for
\texttt{olmo-2-1124-7b-instruct} and \texttt{-13b-instruct}, executed through the
normal HELM path and producing standard artifacts rather than probe output. They
support the attribution's direction without closing it: float32 moves the strict
accuracy from $0.7911$ to $0.7597$ against an official $0.6929$ on the 7B, and
from $0.8555$ to $0.8121$ against $0.7298$ on the 13B --- about a third of the
drift in each case. Precision is therefore a real contributor to the
\texttt{ifeval} gap and not the whole of it, which is a sharper claim than
\S\ref{sec:olmo}(c) makes and a weaker one than a completed confirmation would.

They do not discharge the protocol. Both are \texttt{ifeval} only, on the full
instance set rather than held-out instances, with no second benchmark; neither
has been through the comparison pipeline, so no pair report exists; and the
container ran by local image id rather than a pinned digest, so the runs are not
reproducible across machines. We therefore leave the confirmation open and do not
promote (c) beyond a discovery.
{\color{orange}\textsf{\small[\textsc{pending}: analyse the two float32 runs; extend to held-out
instances, a second benchmark, and a digest-pinned image]}}

The full protocol
separates discovery from held-out confirmation: freeze the discovery instances
and the winning candidate configuration; determine whether the win requires the
probe-only \texttt{add\_special\_tokens=false} behavior, and if so implement it
through the ordinary tokenizer and client path; execute the exact public
\texttt{run\_spec.json} through the normal HELM path rather than the probe,
generating standard artifacts; compare through the standard pair-report
pipeline; evaluate on \emph{held-out} instances and, preferably, a second
benchmark; and preserve the full candidate surface rather than only the winner.
A success would support a concrete attribution-and-closure claim; a failure would
itself be publishable, since it would show that output-matching probes do not
automatically translate into faithful benchmark execution.
\par\smallskip\noindent\fcolorbox{orange}{yellow!12}{%
  \begin{minipage}{0.97\linewidth}\small\textsf{\textbf{[SCAFFOLD --- to be completed]}}\ Insert the confirmation results table: for one dense and one MoE model
$\times$ two benchmarks, report \{naive, exact-recipe, reconstructed-deployment\}
agreement on held-out instances, with exact-token, normalized-text, first-divergent-token,
common-prefix-length, and task-metric columns. {\color{orange}\textsf{\small[\textsc{pending}: held-out confirmation run]}}\end{minipage}}\par\smallskip

\paragraph{Store freshness and preservation.} Some result stores predate the code
that should have produced them: the OLMo aggregate report is stale, and the
GPT-OSS store is a candidate pending a manifest-based acceptance audit. The raw
run artifacts survive in both cases --- every run path recorded in the stored
comparison manifests resolves on disk --- so these are re-render-and-hash tasks
rather than re-executions. What has not been established is that the surviving
artifacts are byte-identical to the ones that produced the stored reports. Until
that audit is done we report these as candidates, not results, and cite no number
that rests on an unaudited store. {\color{orange}\textsf{\small[\textsc{pending}: re-render, hash, and audit the OLMo and
GPT-OSS stores]}}

\paragraph{The audit tool is itself a source of disagreement.} Every number here
passes through our own comparison layer, which decides which stored rows
constitute a run, how completions are normalized before they are compared, and at
what granularity metric keys are matched. Those decisions have been wrong: an
un-deduplicated local join once halved an aggregate that a correct join
reproduces. They are not a coordinate of the original run --- they are a threat to
the validity of our measurement of it --- so we control them by building the
comparison core additively and proving it byte-equal to the incumbent on fixtures
before switching any renderer to it (\S\ref{sec:diff}), and we report a
disagreement only where that core and its predecessor agree.

\paragraph{The exclusion profile is not reported.} \S\ref{sec:preconditions} makes
a typed exclusion a first-class outcome of the taxonomy, but this draft does not
populate it: no corpus-wide breakdown of which runs were excluded, and for which
reason, appears in \S\ref{sec:cases}, and the case studies state only their own
exclusions. The machinery exists and is what produced the scoped corpora used
here; what is missing is a pass over a frozen manifest whose output is fit to
publish. A headline ``fraction reproducible'' statistic would require that
manifest with an explicit numerator and denominator, and fresh stores across
families, so we make no population-level prevalence claim.
{\color{orange}\textsf{\small[\textsc{pending}: corpus-wide exclusion-reason breakdown from a pinned manifest]}}

\paragraph{Single-machine measurements.} Reported agreement is
official-vs-local and same-machine; cross-machine variance (which would separate
platform-specific from recipe-specific drift) is scoped but not completed.
{\color{orange}\textsf{\small[\textsc{pending}: cross-machine repeatability baseline]}}

\paragraph{Oracle objective without log-probs.} Because the hosted originals store
no log-probabilities, the deployment-search objective is completion-text agreement
plus a collapse detector, not a likelihood. This is adequate for ranking
candidates but is not a probabilistic identification.

\paragraph{Non-identifiability is evidence-relative.} The classic-model result
(\S\ref{sec:classic}) is non-identifiability \emph{from the surviving public and
repository evidence we examined}; it does not preclude recovery from a private
provider or author archive we do not have access to.

\paragraph{Score validity is out of scope.} We ask whether a published number can
be reconstructed, not whether it measures what it claims. Open corpora invite
contamination, and a benchmark overstates ability whenever its test data has
reached training, which is rarely checked per benchmark
\citep{sainz2023contamination}. That is an independent axis --- a contaminated
benchmark can be perfectly reproducible and an uncontaminated one irreproducible
--- so we treat the released numbers as the historical record of what HELM
reported and make no claim about their construct validity.

\paragraph{Attribution breadth.} The confirmed attributions are on OLMo-family
\texttt{ifeval}/classic scenarios; ``precision is \emph{the} dominant hidden
variable'' is a hypothesis supported within this pilot, not established across
benchmarks or model families. The efficient-benchmarking thread
(\S\ref{sec:related}) is out of scope here and reserved for future work.
